\subsection{Divisors of zero in a polynomial ring}

PROPOSITION 7. --- Let \(f \in A[X]\) be a non-zero polynomial in an indeterminate and \(\alpha\) its leading coefficient. If \(\alpha\) is cancellable in \(A\) (in particular if \(f\) is monic) then we have for any non-zero element \(g\) of \(A[X]\) ,

\[
f g \neq 0 \quad a n d \quad \deg (f g) = \deg f + \deg g.
\]

Let \(g \in A[X]\) be a non-zero polynomial, \(\beta\) its leading coefficient, \(n = \deg f\) and \(p = \deg g\) . Then the coefficient of \(X^{n+p}\) in \(fg\) is \(\alpha\beta\) which does not vanish, whence the proposition.

PROPOSITION 8. -- If A is an integral domain, then so is \(\mathbf{A}[(\mathbf{X}_i)_{i\in I}]\) .

Let \(u, v\) be two non-zero elements of \(\mathbf{A}[(\mathbf{X}_i)_{i \in I}]\) ; we have to show that \(uv \neq 0\) . Now \(u\) and \(v\) belong to a ring \(\mathbf{A}[(\mathbf{X}_j)_{j \in J}]\) where \(J\) is a finite subset of \(I\) . Thus we can limit ourselves to the case when \(I\) is finite and equal to \(\{1, 2, \ldots, p\}\) . On the other hand, the ring \(\mathbf{A}[\mathbf{X}_1, \ldots, \mathbf{X}_p]\) is isomorphic to the polynomial ring in \(\mathbf{X}_p\) with coefficients in \(\mathbf{A}[\mathbf{X}_1, \ldots, \mathbf{X}_{p-1}]\) . By induction on \(p\) we are thus reduced to proving the proposition for \(\mathbf{A}[\mathbf{X}]\) , and now it suffices to apply Prop. 7.

COROLLARY 1. --- If A is an integral domain, and u, v are elements of \(\mathbf{A}[(\mathbf{X}_{i})_{i\in I}]\) , then \(\deg(uv)=\deg u+\deg v\) .

We can limit ourselves to the case where \(u\) and \(v\) are non-zero. Let \(m = \deg u\) , \(n = \deg v\) ; we have

\[
u = u _ {0} + u _ {1} + \dots + u _ {m}, v = v _ {0} + v _ {1} + \dots + v _ {n}
\]

where \(u_{h}\) (resp. \(v_{h}\) ) is the homogeneous component of degree h of u (resp. v). Since \(u_{m} \neq 0\) and \(v_{n} \neq 0\) , we have \(u_{m}v_{n} \neq 0\) (Prop. 8). Now \(uv = u_{m}v_{n} + w\) with \(\deg w < m + n\) , whence the result.

COROLLARY 2. --- If A is an integral domain, the invertible elements of A{[}(X \(_{i}\) ) \(_{i \in I}\) {]} are the invertible elements of A.

This follows immediately from Cor. 1.

PROPOSITION 9. --- Let \(u \in A[(X_i)_{i \in I}]\) ; then for \(u\) to be nilpotent in the ring \(A[(X_i)_{i \in I}]\) it is necessary and sufficient for all its coefficients to be nilpotent in the ring \(A\) .

As in the proof of Prop. 8 we can make a reduction to the case of polynomials in one variable X. If all the coefficients of u are nilpotent, then u is nilpotent (I, p.~99, Cor. 1). Suppose that \(u\) is nilpotent but not zero and let \(n\) be its degree; we shall argue by induction on \(n\) . Let \(\alpha\) be the leading coefficient of \(u\) . There exists an integer \(m > 0\) such that \(u^m = 0\) . The leading coefficient of \(u^m\) is \(\alpha^m\) , hence \(\alpha^m = 0\) . Now \(u - \alpha X^n\) is nilpotent (I, loc. cit.) and the induction hypothesis shows all the coefficients of \(u - \alpha X^n\) to be nilpotent. Thus all the coefficients of \(u\) are nilpotent.

Remark. --- Let u and v be elements of \(\mathbf{A}[(\mathbf{X}_{i})_{i\in I}]\) , and suppose that A is an integral domain, v is a non-zero multiple of u and v is homogeneous; then u is also homogeneous. For let \(u'\in\mathbf{A}[(\mathbf{X}_{i})_{i\in I}]\) be such that \(v=uu'\) ; we have \(u\neq0\) , \(u'\neq0\) , and if

\[
\begin{array}{c} u = u _ {h} + u _ {h + 1} + \dots + u _ {k} \\ u ^ {\prime} = u _ {h ^ {\prime}} ^ {\prime} + u _ {h ^ {\prime} + 1} ^ {\prime} + \dots + u _ {k ^ {\prime}} ^ {\prime} \end{array}
\]

are the decompositions of u and \(u'\) into homogeneous components, with \(u_{h} \neq 0\) , \(u_{k} \neq 0\) , \(u_{h'}' \neq 0\) , \(u_{k'}' \neq 0\) , then \(v = u_{h} u_{h'}' + u_{h} u_{h'+1}' + \cdots + u_{k} u_{k'}'\) and \(u_{h} u_{h'}'\) is non-zero homogeneous of degree \(h + h'\) while \(u_{k} u_{k'}'\) is non-zero homogeneous of degree \(k + k'\) (Prop. 8). Since v is homogeneous, we have \(h + h' = k + k'\) whence h = k, \(h' = k'\) .

